\section{Visualizations}
\label{sec:proposed-updates}

Visualizing optimizers is inherently difficult because their proposed updates
are functions of the full optimization trajectory.  In this section we try to
peek into the decisions made by the LSTM optimizer, trained on the neural art
task.

\paragraph{Histories of updates}

We select a single optimizee parameter (one color channel of one pixel in the
styled image) and trace the updates proposed to this coordinate by the LSTM
optimizer over a single trajectory of optimization.  We also record the updates
that would have been proposed by both SGD and ADAM if they followed the same
trajectory of iterates.  Figure~\ref{fig:analysis-a} shows the trajectory of
updates for two different optimizee parameters.  From the plots it is clear
that the trained optimizer makes bigger updates than SGD and ADAM.  It is also
visible that it uses some kind of momentum, but its updates are more noisy than
those proposed by ADAM which may be interpreted as having a shorter time-scale
momentum.

\begin{figure}
\centering
\includegraphics[width=\linewidth]{plots/neural-art-updates.pdf}
\caption{Updates proposed by different optimizers at different optimization
steps for two different coordinates.}
\label{fig:analysis-a}
\end{figure}

\paragraph{Proposed update as a function of current gradient}

Another way to visualize the optimizer behavior is to look at the proposed
update $g_t$ for a single coordinate as a function of the current gradient
evaluation $\nabla_t$.  We follow the same procedure as in the previous
experiment, and visualize the proposed updates for a few selected time steps.

These results are shown in
Figures~\ref{fig:coordinate-1}--\ref{fig:coordinate-3}.  In these plots, the
$x$-axis is the current value of the gradient for the chosen coordinate, and
the $y$-axis shows the update that each optimizer would propose should the
corresponding gradient value be observed.  The history of gradient observations
is the same for all methods and follows the trajectory of the LSTM optimizer.

The shape of this function for the LSTM optimizer is often step-like, which is
also the case for ADAM.  Surprisingly the step is sometimes in the opposite
direction as for ADAM, i.e.\ the bigger the gradient, the bigger the update.

\begin{figure}
\centering
\includegraphics[width=\linewidth]{plots/coordinate-1.pdf}
\caption{The proposed update direction for a single coordinate over 32 steps.}
\label{fig:coordinate-1}
\end{figure}

\begin{figure}
\centering
\includegraphics[width=\linewidth]{plots/coordinate-2.pdf}
\caption{The proposed update direction for a single coordinate over 32 steps.}
\label{fig:coordinate-2}
\end{figure}

\begin{figure}
\centering
\includegraphics[width=\linewidth]{plots/coordinate-3.pdf}
\caption{The proposed update direction for a single coordinate over 32 steps.}
\label{fig:coordinate-3}
\end{figure}

